\begin{blocksection}
\question Complete the definition for righttriangle, which takes in a number n and recursively prints out the levels of a right triangle of height n. 

\begin{lstlisting}
def right_triangle(n):
    '''Takes in a non-negative integer n and recursively outputs the layers of a pyramid. The function does not need to return anything. 

    >>> right_triangle(0)
    
    >>> right_triangle(3)
    *
    **
    ***
    >>> num_digits(6)
    *
    ** 
    ***
    **** 
    ***** 
    ****** 
    '''
\end{lstlisting}

\begin{solution}[1in]
\begin{lstlisting}:
    def right_triangle(n):
        if n == 0:
            return
        right_triangle(n-1)
        print('*' * n)
\end{lstlisting}
\end{solution}
\end{blocksection}


\begin{questionmeta}
    \begin{itemize}
        \item This question is meant to help students develop a basic visual intuition behind recursion and a super good question to walk through problem cases with
        \item Questions to ask your students:
        \subitem 1. Why do we print after calling righttriangle?
        \subitem 2. How could we instead print out a mirrored version of the right triangle? 
        \subitem 3. Why aren't we returning anything? 
    \end{itemize}
\end{questionmeta}

